% IAC 2026 manuscript skeleton — master LaTeX source.
% Style: article + iac2026.sty (Word template is visual reference only; no official IAF .cls).
% Scope / cut-keep: SCOPE.md
\documentclass[10pt,twocolumn,letterpaper]{article}
\usepackage{iac2026}

% --- edit after paper-code allocation ---
\renewcommand{\IACpapercode}{IAC-26,A3,IP,xxx,xXXXXX}

\begin{document}

\IACmaketitle
{ARTPS: Operationally Safe Target Prioritization via Depth-Enhanced Hybrid Anomaly Detection}
{Poyraz Baydemir}
{Sel\c{c}uk University, Konya, T\"urkiye}
{%
Planetary rovers operate under strict bandwidth and energy constraints with delayed ground
feedback, which motivates reliable onboard prioritization of scientific targets. Visual-only
autonomy is brittle in field imagery: illumination changes, cast shadows, and low-texture
regions can trigger false positives, while diversity heuristics can suppress high-value
anomalies that appear similar to previously selected targets.
This paper presents ARTPS, an edge-oriented framework that couples monocular relative-depth
estimation with multi-cue anomaly detection and robust localization.
ARTPS fuses image and depth cues via entropy-weighted dynamic fusion, ranks candidates with
a curiosity score and soft similarity penalty, and uses a Priority Buffer for second-chance
re-evaluation of high-anomaly targets down-weighted by diversity policy.
Operational gates address rover-body and boundary-shadow false positives, protect objects
inside shadows, and apply an image-relative size--distance policy based on apparent size
(not metric distance).
On NASA Mars rover imagery, ARTPS reports 0.894 AUROC, 0.847 AUPRC, and 0.823 F1 versus a
PaDiM/PatchCore baseline (0.856 AUROC).
In a lightweight configuration excluding learned depth and autoencoder inference, the core
pipeline sustains 28.1~FPS at $256{\times}256$ on a development workstation.
Proxy ablations (not human bounding-box ground truth) indicate reduced shadow-dense false
positives and lower field-scale false-positive rate under the size--distance policy.
}%
{Mars rover; autonomous exploration; anomaly detection; relative depth; target prioritization; operational safety}

% =====================================================================
\section{Introduction}
% =====================================================================
% TODO: compress Full_Baydemir_ARTPS.pdf §1--§2 here (Related Work folded in).
% Contributions (draft bullets for rewrite):
% (i) depth-guided enhancement + entropy-weighted multi-cue fusion;
% (ii) curiosity ranking with soft similarity / diversity and Priority Buffer;
% (iii) operational FP gates (rover/shadow) and size--distance policy;
% (iv) edge-oriented profiles with measured workstation throughput.
Planetary surface operations motivate onboard screening of scientifically interesting targets
under communications delay~\citep{estlin2012}.
Standard anomaly detectors provide useful baselines but do not by themselves encode
operational false-positive control, diversity-aware ranking, or relative-depth guidance for
field imagery~\citep{defard2020padim,roth2022patchcore}.
Monocular depth networks supply within-image relative structure rather than calibrated
metres~\citep{ranftl2021dpt}; ARTPS treats depth as relative ordering and apparent-size
features accordingly (see also Experimental Protocol).

% =====================================================================
\section{Material and Methods}
% =====================================================================
\subsection{System architecture}
% Pipeline: RGB → relative depth → depth-guided enhance → cue maps → entropy fusion
% → localization / merge → curiosity + diversity → Priority Buffer.
ARTPS processes an RGB frame through relative-depth estimation, optional enhancement,
multi-cue anomaly map extraction, entropy-weighted fusion, localization with box merging,
and ranked target selection with a Priority Buffer for deferred candidates.

\subsection{Depth-guided enhancement}
% KEEP equation: depth-guided enhancement (compress prose from full §3.2.1).
% Optional Real-ESRGAN: one paragraph only; no Results claim without ablation.
Depth-guided enhancement modulates local contrast using the relative-depth map so that
near-field structure is emphasized without claiming metric range.
As an optional implementation path, Real-ESRGAN may be applied as a profile-aware enhance
step; it is not evaluated as a Methods/Results ablation in this manuscript.

\subsection{Anomaly cues and entropy-weighted fusion}
% KEEP: entropy-weighted fusion.
Image cues (texture/edge, photometric), depth cues (discontinuities, proximity weighting),
and optional learned anomaly maps are combined with entropy-weighted dynamic fusion that
up-weights confident components at runtime~\citep{defard2020padim,roth2022patchcore}.

\subsection{Localization, false-positive gates, and size--distance policy}
% KEEP: size-distance policy; tight merge; rover/shadow FP; object-in-shadow gate.
Localization extracts candidate boxes with tightened merge rules.
Rover-body and boundary-shadow masks suppress common false positives; an object-in-shadow
gate protects shadowed rock-like support from over-suppression.
A size--distance policy uses image-relative \emph{relative far} and \emph{apparent size}
bands (not metric size or calibrated distance) to gate field-scale clutter~\citep{ranftl2021dpt}.

\subsection{Curiosity score, diversity, and Priority Buffer}
% KEEP: curiosity score; soft similarity penalty; uncertainty/disagreement; Priority Buffer.
Candidates are ranked by a curiosity score combining known-class and anomaly terms, with a
soft similarity penalty to reduce redundant selections and optional uncertainty/disagreement
signals.
The Priority Buffer retains high raw-anomaly candidates that diversity policy would otherwise
drop, for second-chance review when resources allow.

% =====================================================================
\section{Experimental Protocol}
% =====================================================================
\subsection{Data and splits}
% Dataset / labeling / split from full §4.1 — rewrite, do not dump.
Evaluation uses NASA Mars rover imagery; labeling and split details will be summarized from
the full manuscript without installation or UI narrative.

\subsection{Baselines, metrics, and hardware}
Baselines include PaDiM/PatchCore-style anomaly detection~\citep{defard2020padim,roth2022patchcore}.
We report AUROC, AUPRC, and F1 by name (no textbook formula expansions).
Hardware timings refer to a development workstation at fixed input resolutions.

\subsection{Ablation and proxy protocol}
% Proxy disclaimer required in Results captions.
Component ablations follow the full-manuscript protocol.
Additional \emph{proxy} ablations for shadow/FP gates and size--distance policy use class
labels, depth/shadow masks, and OFF-run pseudo-ground-truth; they are \textbf{not} human
bounding-box ground truth.
The synthetic size/distance lite bench is software verification only and is not a performance
result.

% =====================================================================
\section{Results}
% =====================================================================
\subsection{Detection, ranking, and localization}
% Fill from full §6.1--6.2; numbers below from portal abstract / prior runs.
On the primary Mars-imagery evaluation, ARTPS reports 0.894 AUROC, 0.847 AUPRC, and
0.823 F1, compared with a PaDiM/PatchCore baseline of 0.856 AUROC.

\subsection{Field conditions and hardware}
% Full §6.3--6.4.
In a lightweight configuration (OpenCV preprocessing + fusion + localization, excluding
learned depth/autoencoder inference), the core pipeline sustains 28.1~FPS at $256{\times}256$
on a development workstation (profile-dependent at higher resolutions).

\subsection{Shadow / FP proxy ablation}
% Wired to results/paper_figs/iac_shadow_proxy_*
\begin{table}[t]
\centering
\caption{Shadow / false-positive proxy ablation ($n{=}21$ images, 6 shadow-dense).
Proxy metrics (class labels, depth/shadow masks, OFF-run pseudo-GT), not human bbox GT.
Depth cues are relative depth ordering within each image.}
\label{tab:shadow-proxy}
\small
\begin{tabular}{@{}lccc@{}}
\toprule
Metric & OFF & ON & $\Delta$ \\
\midrule
Shadow-dense FPR & 0.083 & 0.000 & $-0.083$ \\
Avg.\ detections & 3.48 & 2.76 & --- \\
Target recall (proxy, ON) & --- & 0.313 & --- \\
Shadow-rock loss (ON) & --- & 0.000 & --- \\
\bottomrule
\end{tabular}
\end{table}
Source: \path{results/paper_figs/iac_shadow_proxy_summary.json}.

\subsection{Size--distance policy proxy ablation}
% Wired to results/paper_figs/iac_size_distance_proxy_*
\begin{table}[t]
\centering
\caption{Size--distance policy proxy ablation ($n{=}21$).
Bands are image-relative; features use apparent size (not metric size).
Lite bench is software verification only.}
\label{tab:size-distance-proxy}
\small
\begin{tabular}{@{}lccc@{}}
\toprule
Metric & Policy OFF & Policy ON & $\Delta$ \\
\midrule
Field-scale FPR & 0.500 & 0.351 & $-0.149$ \\
Near-large over-merge & 0.000 & 0.000 & $0.000$ \\
Mean matched IoU (self) & --- & 0.999 & --- \\
Avg.\ detections & 2.29 & 2.76 & --- \\
\bottomrule
\end{tabular}
\end{table}
Source: \path{results/paper_figs/iac_size_distance_proxy_summary.json}.

% =====================================================================
\section{Discussion}
% =====================================================================
% Fold full §7 Limitations + §8 Safety / flight-readiness here.
Proxy improvements in shadow-dense FPR and field-scale FPR support the operational gates,
but proxy labels are not a substitute for human bbox evaluation.
Relative depth remains within-image ordering; metric range and cross-image depth comparison
are out of scope until scale-aware calibration exists.
Priority Buffer and diversity trade false-negative risk against redundancy under bandwidth
limits; failure cases (textureless terrain, severe illumination) should be enumerated in the
full rewrite.
Flight-readiness claims stay limited to profile-aware workstation screening, not flight
hardware certification.

% =====================================================================
\section{Conclusions}
% =====================================================================
% Short; no industrial applications laundry list.
ARTPS couples relative-depth cues, entropy-weighted fusion, curiosity ranking, and
operational buffers/gates for safer onboard target prioritization on Mars rover imagery.
Future work includes human-annotated localization benchmarks and scale-aware depth if
mission calibration becomes available.

\bibliographystyle{unsrtnat}
\bibliography{refs}

\end{document}
